% GENERATED FILE. Edit canonical lessons, then run npm run generate:books.
% canonical-source-hash: fnv1a64:7764aab169c686f6
% canonical-lessons: JA-C250-mita, JA-C250-katta, JA-R250-what-i-did-yesterday

\chapter{What I Did Yesterday}
\label{ch:ja-what-i-did-yesterday}
\hlchaptermodality{250}

\begin{chapteropening}
\textbf{By the end of this chapter:} \emph{I can say what I did, to a friend, with the plain past: mita, tabeta, shita, kita, and katta, nonda, kaita, itta.}

Complete JA-R250-what-i-did-yesterday: tell a friend what you did yesterday, in the plain past.
\end{chapteropening}

\section[\texorpdfstring{\ja{みた} (mita)}{みた (mita)}]{\ja{みた} (mita) — the plain past, what friends say}

\label{lesson:JA-C250-mita}



\begin{quote}
Before the new form: say \textquotedblleft{}I saw a film yesterday\textquotedblright{}, politely. (\emph{kinō eiga o mimashita}.) The polite past is what you say to a teacher or in a shop. To a friend, the past is shorter.
\end{quote}

\subsection*{You'll want to know: the plain past}
A friend asks what you did yesterday, and you answer as friends do:

\begin{quote}
\textbf{\ja{えいがを} \ja{みた。}}
\end{quote}

\begin{quote}
\emph{eiga o mita.} — \textquotedblleft{}I saw a film.\textquotedblright{}
\end{quote}

\begin{quote}
\textbf{\ja{きのう} \ja{えいがを} \ja{みた。}}
\end{quote}

\begin{quote}
\emph{kinō eiga o mita.} — \textquotedblleft{}I saw a film yesterday.\textquotedblright{}
\end{quote}

\textbf{\ja{みた}} (\emph{mita}) is \textbf{\ja{みました}} (\emph{mimashita}) without the politeness: the plain past, the form friends and family use with each other.

For \textbf{\ja{みる}} (\emph{miru}), \textquotedblleft{}to see\textquotedblright{}, drop the last \textbf{\ja{る}} and add \textbf{\ja{た}}. \textbf{\ja{たべる}} (\emph{taberu}), \textquotedblleft{}to eat\textquotedblright{}, goes the same way: \textbf{\ja{たべた}} (\emph{tabeta}), \textquotedblleft{}ate\textquotedblright{}.

\begin{quote}
\textbf{\ja{ばんごはんを} \ja{たべた。}}
\end{quote}

\begin{quote}
\emph{bangohan o tabeta.} — \textquotedblleft{}I ate dinner.\textquotedblright{}
\end{quote}

\begin{quote}
\textbf{\ja{うちで} \ja{ばんごはんを} \ja{たべた。}}
\end{quote}

\begin{quote}
\emph{uchi de bangohan o tabeta.} — \textquotedblleft{}I ate dinner at home.\textquotedblright{}
\end{quote}

The polite form shows the way too: take \textbf{\ja{ます}} off \textbf{\ja{ねます}} (\emph{nemasu}) and add \textbf{\ja{た}}: \textbf{\ja{ねた}} (\emph{neta}), \textquotedblleft{}went to bed\textquotedblright{}. From \textbf{\ja{おきます}} (\emph{okimasu}) comes \textbf{\ja{おきた}} (\emph{okita}), \textquotedblleft{}got up\textquotedblright{}.

\begin{quote}
\textbf{\ja{きのう} \ja{はやく} \ja{ねた。}}
\end{quote}

\begin{quote}
\emph{kinō hayaku neta.} — \textquotedblleft{}I went to bed early yesterday.\textquotedblright{}
\end{quote}

\subsection*{You'll want to know: \ja{した} and \ja{きた}}
Two more verbs: \textbf{\ja{した}} (\emph{shita}), \textquotedblleft{}did\textquotedblright{}. It is the past of \textbf{\ja{する}} (\emph{suru}), \textquotedblleft{}to do\textquotedblright{}, the verb inside \textbf{\ja{します}} (\emph{shimasu}). And \textbf{\ja{くる}} (\emph{kuru}), \textquotedblleft{}to come\textquotedblright{}, the verb inside \textbf{\ja{きます}} (\emph{kimasu}), has \textbf{\ja{きた}} (\emph{kita}), \textquotedblleft{}came\textquotedblright{}.

\begin{quote}
\textbf{\ja{しゅくだいを} \ja{した。}}
\end{quote}

\begin{quote}
\emph{shukudai o shita.} — \textquotedblleft{}I did my homework.\textquotedblright{}
\end{quote}

\begin{quote}
\textbf{\ja{ともだちが} \ja{きた。}}
\end{quote}

\begin{quote}
\emph{tomodachi ga kita.} — \textquotedblleft{}A friend came.\textquotedblright{}
\end{quote}

\begin{quote}
\textbf{\ja{きのう} \ja{ともだちが} \ja{きた。}}
\end{quote}

\begin{quote}
\emph{kinō tomodachi ga kita.} — \textquotedblleft{}A friend came yesterday.\textquotedblright{}
\end{quote}

The plain past has a second job, and later lessons give it to you: it goes inside a polite sentence, in front of the word that makes the sentence polite. So \textbf{\ja{みた}} is worth owning even if you speak politely to everyone.

\subsection*{Your turn}
Say these aloud:

\begin{itemize}
\raggedright
  \item \textquotedblleft{}I ate dinner\textquotedblright{} — \emph{bangohan o tabeta}
  \item \textquotedblleft{}I ate breakfast\textquotedblright{} — \emph{asagohan o tabeta}
  \item \textquotedblleft{}I ate lunch\textquotedblright{} — \emph{hirugohan o tabeta}
\end{itemize}

\begin{tcolorbox}[breakable,colback=teal!4,colframe=teal!35!black,title={Before you move on}]
\textquotedblleft{}I saw\textquotedblright{}, to a friend? (\textbf{\ja{みた}}, \emph{mita}.) \textquotedblleft{}I got up\textquotedblright{}? (\textbf{\ja{おきた}}, \emph{okita}.) \textquotedblleft{}Did\textquotedblright{}? (\textbf{\ja{した}}, \emph{shita}.) \textquotedblleft{}Came\textquotedblright{}? (\textbf{\ja{きた}}, \emph{kita}.) And \textquotedblleft{}I saw\textquotedblright{}, to a teacher? (\textbf{\ja{みました}}, \emph{mimashita}.)
\end{tcolorbox}
\section[\texorpdfstring{\ja{かった} / \ja{のんだ} … (katta / nonda …)}{かった / のんだ … (katta / nonda …)}]{\ja{かった} / \ja{のんだ} / \ja{かいた} (katta / nonda / kaita) — the past by the last kana}

\label{lesson:JA-C250-katta}



\begin{quote}
Say \textquotedblleft{}I ate\textquotedblright{} and \textquotedblleft{}I did\textquotedblright{}, to a friend. (\emph{tabeta}, \emph{shita}.) Then say \textquotedblleft{}I bought shoes\textquotedblright{}, politely. (\emph{kutsu o kaimashita}.) Most verbs do not drop a \textbf{\ja{る}}: their \textbf{-\ja{た}} depends on their last kana.
\end{quote}

\subsection*{You'll want to know: \ja{った} and \ja{んだ}}
\textbf{\ja{かう}} (\emph{kau}), \textquotedblleft{}to buy\textquotedblright{}, ends in \textbf{\ja{う}}. Its plain past is \textbf{\ja{かった}} (\emph{katta}): a small \textbf{\ja{っ}}, a short catch, then \textbf{\ja{た}}.

\begin{quote}
\textbf{\ja{くつを} \ja{かった。}}
\end{quote}

\begin{quote}
\emph{kutsu o katta.} — \textquotedblleft{}I bought shoes.\textquotedblright{}
\end{quote}

\textbf{\ja{つ}} and \textbf{\ja{る}} go the same way. \textbf{\ja{まつ}} (\emph{matsu}), \textquotedblleft{}to wait\textquotedblright{}, has \textbf{\ja{まった}} (\emph{matta}). \textbf{\ja{かえります}} (\emph{kaerimasu}) is the polite form of \textbf{\ja{かえる}} (\emph{kaeru}), \textquotedblleft{}to go home\textquotedblright{}, and its past is \textbf{\ja{かえった}} (\emph{kaetta}). (A different \ja{かえる}, \textquotedblleft{}to change\textquotedblright{}, drops its \ja{る} like \textbf{\ja{みる}}.)

\begin{quote}
\textbf{\ja{えきで} \ja{まった。}}
\end{quote}

\begin{quote}
\emph{eki de matta.} — \textquotedblleft{}I waited at the station.\textquotedblright{}
\end{quote}

\begin{quote}
\textbf{\ja{はやく} \ja{うちに} \ja{かえった。}}
\end{quote}

\begin{quote}
\emph{hayaku uchi ni kaetta.} — \textquotedblleft{}I went home early.\textquotedblright{}
\end{quote}

\textbf{\ja{む}} and \textbf{\ja{ぶ}}, and the rare verbs in \textbf{\ja{ぬ}}, end in \textbf{\ja{んだ}}. \textbf{\ja{のむ}} (\emph{nomu}), \textquotedblleft{}to drink\textquotedblright{}, has \textbf{\ja{のんだ}} (\emph{nonda}); \textbf{\ja{あそぶ}} (\emph{asobu}), \textquotedblleft{}to play\textquotedblright{}, has \textbf{\ja{あそんだ}} (\emph{asonda}).

\begin{quote}
\textbf{\ja{おちゃを} \ja{のんだ。}}
\end{quote}

\begin{quote}
\emph{ocha o nonda.} — \textquotedblleft{}I drank tea.\textquotedblright{}
\end{quote}

\begin{quote}
\textbf{\ja{こうえんで} \ja{あそんだ。}}
\end{quote}

\begin{quote}
\emph{kōen de asonda.} — \textquotedblleft{}I played in the park.\textquotedblright{}
\end{quote}

\subsection*{You'll want to know: \ja{いた}, \ja{いだ}, \ja{した}, and \ja{いった}}
A verb in \textbf{\ja{く}} ends in \textbf{\ja{いた}}: \textbf{\ja{かく}} (\emph{kaku}), \textquotedblleft{}to write\textquotedblright{}, has \textbf{\ja{かいた}} (\emph{kaita}); \textbf{\ja{きく}} (\emph{kiku}), \textquotedblleft{}to listen\textquotedblright{}, has \textbf{\ja{きいた}} (\emph{kiita}). One in \textbf{\ja{ぐ}} ends in \textbf{\ja{いだ}}: \textbf{\ja{およぐ}} (\emph{oyogu}), \textquotedblleft{}to swim\textquotedblright{}, has \textbf{\ja{およいだ}} (\emph{oyoida}).

\begin{quote}
\textbf{\ja{てがみを} \ja{かいた。}}
\end{quote}

\begin{quote}
\emph{tegami o kaita.} — \textquotedblleft{}I wrote a letter.\textquotedblright{}
\end{quote}

\begin{quote}
\textbf{\ja{うみで} \ja{およいだ。}}
\end{quote}

\begin{quote}
\emph{umi de oyoida.} — \textquotedblleft{}I swam in the sea.\textquotedblright{}
\end{quote}

A verb in \textbf{\ja{す}} ends in \textbf{\ja{した}}: \textbf{\ja{はなす}} (\emph{hanasu}), \textquotedblleft{}to speak\textquotedblright{}, has \textbf{\ja{はなした}} (\emph{hanashita}).

\begin{quote}
\textbf{\ja{ともだちと} \ja{はなした。}}
\end{quote}

\begin{quote}
\emph{tomodachi to hanashita.} — \textquotedblleft{}I talked with a friend.\textquotedblright{}
\end{quote}

And one verb in \textbf{\ja{く}} goes its own way: \textbf{\ja{いく}} (\emph{iku}), \textquotedblleft{}to go\textquotedblright{}, has \textbf{\ja{いった}} (\emph{itta}), like \textbf{\ja{かった}}.

\begin{quote}
\textbf{\ja{きのう} \ja{やまに} \ja{いった。}}
\end{quote}

\begin{quote}
\emph{kinō yama ni itta.} — \textquotedblleft{}I went to the mountains yesterday.\textquotedblright{}
\end{quote}

\subsection*{Your turn}
Say these aloud:

\begin{itemize}
\raggedright
  \item \textquotedblleft{}I went to the sea\textquotedblright{} — \emph{umi ni itta}
  \item \textquotedblleft{}I went to the river\textquotedblright{} — \emph{kawa ni itta}
  \item \textquotedblleft{}I went to the park\textquotedblright{} — \emph{kōen ni itta}
\end{itemize}

\begin{tcolorbox}[breakable,colback=teal!4,colframe=teal!35!black,title={Before you move on}]
\textquotedblleft{}Bought\textquotedblright{}? (\textbf{\ja{かった}}, \emph{katta}.) \textquotedblleft{}Drank\textquotedblright{}? (\textbf{\ja{のんだ}}, \emph{nonda}.) \textquotedblleft{}Wrote\textquotedblright{}? (\textbf{\ja{かいた}}, \emph{kaita}.) \textquotedblleft{}Swam\textquotedblright{}? (\textbf{\ja{およいだ}}, \emph{oyoida}.) \textquotedblleft{}Talked\textquotedblright{}? (\textbf{\ja{はなした}}, \emph{hanashita}.) \textquotedblleft{}Went\textquotedblright{}? (\textbf{\ja{いった}}, \emph{itta}.)
\end{tcolorbox}
\section[(review)]{Review — what I did yesterday}

\label{lesson:JA-R250-what-i-did-yesterday}



\begin{quote}
Say \textquotedblleft{}I saw\textquotedblright{} and \textquotedblleft{}I bought\textquotedblright{}, to a friend. (\emph{mita}, \emph{katta}.)
\end{quote}

\subsection*{Your turn}
Say these aloud:

\begin{itemize}
\raggedright
  \item I saw a film yesterday. — to a friend
  \item I got up early and ate breakfast.
  \item A friend came, and we played in the park.
  \item I drank tea and talked with a friend.
  \item I went to the sea and swam.
\end{itemize}

\begin{tcolorbox}[breakable,colback=teal!4,colframe=teal!35!black,title={Before you move on}]
Two sentences for each of those, or one? (Two: \textbf{\ja{おきた。あさごはんを} \ja{たべた。}} The form that joins them comes later.) \textquotedblleft{}Waited\textquotedblright{}? (\textbf{\ja{まった}}, \emph{matta}.) \textquotedblleft{}Listened\textquotedblright{}? (\textbf{\ja{きいた}}, \emph{kiita}.) \textquotedblleft{}Played\textquotedblright{}? (\textbf{\ja{あそんだ}}, \emph{asonda}.) \textquotedblleft{}Went home\textquotedblright{}? (\textbf{\ja{かえった}}, \emph{kaetta}.)
\end{tcolorbox}
